\section{Relations between types of smoothings}
\label{sec:smoothing-types}

Before proceeding with the main proof, we first examine the relationships between different types of smoothing. In particular, we will consider the distinctions between oriented and unoriented smoothing, as well as between connected and disconnected smoothing, and explore how these two attributes interact.

A priori the hypothesis of oriented smoothing on $f$ is weaker than
smoothing, since it does not include the hypotheses on unoriented smoothings.
However, it follows from the other hypotheses; the key ingredient is
our previous work in~\cite{MGT21:Smoothings}. We prove the following,
which in particular will allow the simpler algebraic statement in
Theorem~\ref{thm:intersections_group}; it is not necessary for the
proof of Theorem~\ref{thm:intersect}.

\begin{proposition}
If $f$ satisfies oriented smoothing
(connected and disconnected), convex union,
stability, homogeneity, and symmetry, then $f$ satisfies smoothing.
\label{prop:oriented_smoothing_unoriented}
\end{proposition}

The proof uses the following result, which can be seen as a weak
converse of the main result in our previous paper.
\begin{proposition}
Suppose $f$ is an symmetric curve functional satisfying
convex union and homogeneity that extends continuously to geodesic currents.
If $f$ satisfies disconnected smoothing (resp.\ connected smoothing) then $f$ satisfies connected smoothing (resp.\ disconnected smoothing).
\label{prop:cont_conn_disc}
\end{proposition}

\begin{proof}[Proof of
  Proposition~\ref{prop:oriented_smoothing_unoriented}, assuming Proposition~\ref{prop:cont_conn_disc}]
  Since $f$ satisfies disconnected oriented smoothing and is
  symmetric, it also
  satisfies disconnected unoriented smoothing. Indeed, let
  $(\ora{a},\ora{b})$ R-cross. Then $(\ora{B},\ora{a})$ R-cross, so by
  disconnected oriented smoothing
\[
f([B][a]) \geq f([Ba]).
\]
Since $f$ is symmetric, we have $f([b][a])=f([B][a]) \geq
f([Ba])$, as desired.

Since $f$ satisfies oriented smoothing, convex union, stability and
homogeneity, it extends continuously and homogeneously
to oriented geodesic currents, by~\cite[Theorem~A]{MGT21:Smoothings}. Hence by
Proposition~\ref{prop:cont_conn_disc}, it satisfies connected
smoothing, and thus all cases of smoothing.
\end{proof}

As an immediate corollary, we have the following.

\begin{corollary}
If $f$ is symmetric and satisfies oriented smoothing, additivity, and
stability, then $f$ satisfies smoothing.
\label{cor:oriented_smoothing_unoriented}
\end{corollary}

We now turn to the proof of Proposition~\ref{prop:cont_conn_disc},
starting with some lemmas. The first lets us approximate
disconnected curves by connected curves and vice versa.
\DPTtodo{Earlier (Eq. \eqref{eq:delta-curve}) we used the notation
  $\delta_{[g]}$ for this same current. Holding off on switching it
  in this section till we decide which one we prefer. I did move the def out of
  the statement of the lemma}
\DMGtodo{I changed it to $\delta$, and added conjugacy class to the equation}
\begin{lemma}
  Given any $g,h \in \pi_1(S)$, we have
\[
\lim_{n\to\infty} \frac{\delta_{[g^n h^n]}}{n} = \delta_{[g]} + \delta_{[h]}
\]
in the weak$^*$-topology of geodesic currents.
\label{lem:gnhn_to_gsumh}
\end{lemma}

Recall that here $\delta_{[g]}$ denotes the geodesic current representing the curve $[g]$ (Equation~\ref{eq:delta-curve}). \DMGtodo{added this reference}

\begin{proof}
  This is a straightforward fact of the weak$^*$-topology convergence
  of geodesic currents and follows from computations similar
  to~\cite[Proposition~5.1]{Sas22:Currents}.
\end{proof}

Next we give some extensions of Lemma~\ref{lem:righthandedfoundation}
to powers.
\begin{lemma}
  Let  $a,b \in \pi_1(S)$ so that $\bigl(\ora{a},\ora{b}\bigr)$ are
  R-parallel. Then for all
  $n \ge 1$, the triple $(\ora{a},\ora{a^n b^n}, \ora{b})$ is R-parallel.
  \label{lem:powers_parallel}
\end{lemma}
\begin{proof}
The proof is by induction on~$n$, repeatedly applying
Lemma~\ref{lem:righthandedfoundation}\ref{item:ab-parallel}. The case
$n=1$ is directly that lemma.
For the inductive step, we have:
\DPTtodo{I started moving the reasons next to the equations. Only did
  a few to see if we like it}
  \DMGtodo{I like it}
\begin{align*}
  &\bigl(\ora{a},\ora{a^{n-1}b^{n-1}},\ora{b}\bigr)\text{ R-parallel}
    &&\text{(Induction)}\\
  &\quad\Rightarrow
    \bigl(\ora{a},\ora{a^n b^{n-1}},\ora{a^{n-1}b^{n-1}},\ora{b}\bigr)\text{ R-parallel}
    &&\text{(Lemma~\ref{lem:righthandedfoundation}\ref{item:ab-parallel} on
$\bigl(\ora{a},\ora{a^{n-1}b^{n-1}}\bigr)$)}\\
  &\quad\Rightarrow
    \bigl(\ora{a},\ora{a^n b^{n-1}},\ora{a^n b^n},\ora{b}\bigr)\text{ R-parallel}
    &&\text{(Lemma~\ref{lem:righthandedfoundation}\ref{item:ab-parallel} on
$\bigl(\ora{a^n b^{n-1}},\ora{b}\bigr)$)}.
\qedhere
\end{align*}
% where we start with the inductive hypothesis, used
% Lemma~\ref{lem:righthandedfoundation}\ref{item:ab-parallel} on
% $\bigl(\ora{a},\ora{a^{n-1}b^{n-1}}\bigr)$, and then forgot about
% $\ora{a^{n-1}b^{n-1}}$ and used 
% Lemma~\ref{lem:righthandedfoundation}\ref{item:ab-parallel} on
% $\bigl(\ora{a^n b^{n-1}},\ora{b}\bigr)$.
\end{proof}

\begin{lemma}
  Let  $a,b \in \pi_1(S)$ so that $\bigl(\ora{a},\ora{b}\bigr)$ are
  R-parallel. Then for all
  $n \ge 1$, $m \geq 1$ the pair $(\ora{b^ma^m},\ora{a^n
  b^n})$ R-cross.
\label{lem:anbn_vs_bkak}
\end{lemma}
\begin{proof}
  Assume first $m=1$.
  We proceed it by induction on~ $n$.
  The base case is
 is clear by
 Lemma~\ref{lem:righthandedfoundation}\ref{item:ab-parallel}. In
 general we have
 \begin{align*}
   &\bigl(\ora{ba},\ora{a^{n-1}b^{n-1}}\bigr)\text{ R-cross}
     &&\text{(Induction)}\\
   &\quad\Rightarrow
     \bigl(\ora{ba},\ora{a^{n-1}b^n a}\bigr)\text{ R-cross}
     &&\text{(Lemma~\ref{lem:rhs_prod})}\\
   &\quad\Rightarrow
     \bigl(\ora{ab}, \ora{a^nb^n}\bigr)\text{ R-cross}
     &&\text{(Conjugate by~$a$)}.
 \end{align*}
% where in the first three steps we start with the inductive hypothesis, use
% Lemma~\ref{lem:rhs_prod}, and
% conjugate by~$a$.
Finally,
Lemma~\ref{lem:powers_parallel} and the R-crossing of 
$\bigl(\ora{ab}, \ora{a^nb^n}\bigr)$ restrict the
endpoints of $\ora{a^n b^n}$ to the indicated intervals in
Figure~\ref{fig:parallel_lemma}, which implies that $\ora{a^n b^n}$ crosses
$\ora{ba}$.

For $m \geq 1$, we apply the symmetric argument to show that 
$\bigl(\ora{b^ma^m}, \ora{ba}\bigr)$ R-cross with the endpoints of $\ora{b^ma^m}$ restricted to the intervals in Figure~\ref{fig:parallel_lemma}.
\end{proof}
\DPTtodo{In figure, labelings of the intervals could be clearer; maybe
  refer to color? $\protect\ora{ab}$ is not labeled. The TeX labels seem to be
  switched. I edited the
  caption to include
  Lemma~\ref{lem:powers_parallel}.}
  \DMGtodo{I labeled $\protect\ora{ab}$ and shifted some labels to make them more visible. I'm not sure I see which ones are switched.}
\begin{figure}
     \begin{subfigure}[b]{0.45\textwidth}
\centerline{\svgfignormal{anbn_ba_c}}
    \caption{$\protect\ora{a}$ and $\protect\ora{b}$ are R-parallel
      implies $(\protect\ora{b^ma^m}, \protect\ora{a^n b^n})$ R-cross
      and $(\ora{a},\ora{a^n b^n}, \ora{b})$ are R-parallel}
    \label{fig:parallel_lemma}
  \end{subfigure}
  \hfill
    \begin{subfigure}[b]{0.45\textwidth}
 \centerline{\svgfignormal{anbn_ba_p}}
    \caption{$\protect\ora{b}$ and $\protect\ora{a}$ R-cross implies $\protect\ora{a^n b^n}$ is R-parallel to $\protect\ora{ba}$}
    \label{fig:cross_lemma}
     \end{subfigure}
  \caption{Crossings of powers}
\end{figure}



\begin{lemma} 
  Let $a,b \in \pi_1(S)$ so that $(\ora{b},\ora{a})$ R-cross. Then for all
  $n \ge 1$, $(\ora{a^n
  b^n},\ora{ba})$ are R-parallel, and for every $n \geq 2$, the pair $(\ora{a^n
  b^n},\ora{ab})$ are R-parallel.
\label{lem:anbn_vs_bkak_parallel}
\end{lemma}
\begin{proof}
We proceed it by induction on $n$. The base case is clear by
Lemma~\ref{lem:righthandedfoundation}\ref{item:ab-cross}.
In general we have
\begin{align*}
  \bigl(\ora{a^{n-1}b^{n-1}},\ora{ba}\bigr)\text{ R-parallel}
  &\Rightarrow
    \bigl(\ora{a^{n-1}b^n a}, \ora{ba}\bigr)\text{ R-parallel}\\
  &\Rightarrow
    \bigl(\ora{a^nb^n}, \ora{ab}\bigr)\text{ R-parallel}\\
  &\Rightarrow
    \bigl(\ora{a^nb^n}, \ora{ba}\bigr)\text{ R-parallel}
\end{align*}
where we use the inductive hypothesis,
Lemma~\ref{lem:righthandedfoundation}\ref{item:ab-parallel},
conjugation, and transitivity of R\hyp parallelism. See
Figure~\ref{fig:cross_lemma}.
\end{proof}


\DPTtodo[inline]{This proof has a slightly complicated structure,
  shown below (out TODO environment to allow a list). I'm wondering if
  there are ways to make this clearer. For instance, maybe have two
  Proof environments.}
  \DMGtodo{I implemented the two proof environment proposal}

\begin{proof}[Proof of Direction 1 of Proposition~\ref{prop:cont_conn_disc}]
		
Suppose first $f \colon \Curves^+(S) \to \mathbb{R}$ satisfies convex union, disconnected smoothing, homogeneity and it extends continuously to geodesic currents.
		\begin{itemize}
			\item \textbf{Part one:} Oriented connected smoothing $[ab] \reducesto
			[a][b]$.
			
			Given a pair $a, b \in \pi_1(S)$ with $(\ora{a},\ora{b})$ parallel (so
			$[ab] \reducesto [a][b]$ and $[ab]\reducesto[aB]$), we show first $f$
			satisfies oriented connected
			smoothing, i.e., $f([ab]) \geq f([a][b])$.
			
			We will start by showing, by induction on $n$, that
			\begin{equation}\label{eq:power_ineq}
				2nf([ab]) \geq f([a^{2n}b^{2n}]).
			\end{equation}
			The base case $n=1$ follows from the fact that if $\ora{a},\ora{b}$
			are parallel, $\ora{ab}$ and $\ora{ba}$ cross, and hence
			\[
			2[ab]=[ab][ba] \reducesto [ab^2a]=[a^2b^2].
			\]
			Hence, by homogeneity and disconnected smoothing, the base case follows.
			Now, we prove the induction step.
			\begin{align*}
				2nf([ab]) &= 2(n-1)f([ab]) + 2f([ab])\\
				&\geq f([a^{2(n-1)}b^{2(n-1)}]) + f([b^2a^2])
				&&\text{(Induction)}\\
				&\geq f([a^{2(n-1)}b^{2(n-1)}][b^2a^2])
				&&\text{(Convex union)}\\
				&\geq f([a^{2(n-1)}b^{2(n-1)} b^2a^2]) = f([a^{2n}b^{2n}])
				&&\text{(Lemma~\ref{lem:anbn_vs_bkak} and disconnected smoothing)}
			\end{align*}
			% In the first inequality we used induction step, in the second convex
			% union, then Lemma~\ref{lem:anbn_vs_bkak} and the assumption that $f$
			% satisfies disconnected smoothing.
			Divide Eq.~\eqref{eq:power_ineq} by $2n$, let $n$ go to infinity, and
			use Lemma~\ref{lem:gnhn_to_gsumh} to obtain
			\[
			f([ab]) \geq f([a] [b])
			\]
			which proves oriented connected smoothing.
			
			\item \textbf{Part two:} Unoriented connected smoothing $[ab] \reducesto
			[aB]$.
			
			
			Second, we prove unoriented connected smoothing: with $a,b$ as above,
			$f([ab]) \ge f([aB])$.
			We split into two cases, depending on the relative position of $\ora{aB}$
			and $\ora{Ba}$. By
			Lemma~\ref{lem:ab_parallel_aB}, $(\ora{aB},\ora{Ba})$ are either anti-parallel
			or cross.
			
			
			\begin{itemize}
				\item \textbf{Case A:} $\ora{aB}, \ora{Ba}$ anti-parallel.
				
				We first consider the case when $(\ora{aB},\ora{Ba})$ are anti-parallel. We will prove by induction that
				\begin{equation}
					2nf([ab]) \geq f([(aB)^n (Ab)^n]).
					\label{eq:antiparallel_induction}
				\end{equation}
				The base case, $n=1$, follows since $(\ora{ab},\ora{ba})$ cross, hence
				\[
				[ba][ab] \reducesto [ba(ab)^{-1}] = [(aB)(Ab)].
				\]
				By homogeneity and disconnected smoothing, we get the base case.
				Now, we prove the induction step:
				\begin{center}
				\begin{align*}
					2nf([ab]) = 2(n-1)f([ab]) + 2f([ab]) &\geq f([(aB)^{n-1} (Ab)^{n-1}]) + f([(Ab)(aB)]) \\
					&\geq f([(aB)^{n-1} (Ab)^{n-1} (Ab)(aB)]) = f([(aB)^n (Ab)^n]).
				\end{align*}
				\end{center}
				In the first inequality we used the induction step. The second inequality
				uses convex union and disconnected smoothing applied to
				\[[(aB)^{n-1}(Ab)^{n-1}][(Ab)(aB)] \reducesto [(aB)^{n-1} (Ab)^{n-1}
				(Ab)(aB)].
				\]
				 For this reduction, note that $\ora{aB}$ and $\ora{Ba}$ are
				anti-parallel, so $\ora{aB}$ and $\ora{Ab}$ are parallel. Hence, we
				can apply Lemma~\ref{lem:anbn_vs_bkak} to conclude that
				$\ora{(aB)^{n-1} (Ab)^{n-1}}$ and $\ora{(Ab)(aB)}$ cross.
				
				Divide Eq.~\eqref{eq:antiparallel_induction} by $2n$, let $n$ go to
				infinity, and use Lemma~\ref{lem:gnhn_to_gsumh} to get
				\[
				f([ab]) \geq \frac{f([aB] [Ab])}{2} = f([aB])
				\]
				where we used symmetry and homogeneity in the last equality.
				
				
				\item \textbf{Case B:} $\ora{aB}, \ora{Ba}$ cross.
				
				Now, still in the unoriented connected smoothing case, suppose
				$(\ora{aB},\ora{Ba})$ cross.
				If $(\ora{a},\ora{b})$ are parallel it follows by Lemma~\ref{lem:a_parallel_Ba} that either $(\ora{a},\ora{bA})$ are parallel or $(\ora{b}, \ora{aB})$ are parallel.
				In the first case, we have 
				\[
				f([ba]) = f([bA)(a^2)]) \geq f([bA]) + f(a^2) \geq f([bA])
				\]
				where in the first inequality we used oriented connected smoothing, and in the last inequality the fact that $f(a^2) \geq 0$, by Proposition~\ref{prop:curve-positive} below.
				In the second case,
				\[
				f([ab]) = f([aB)(b^2)]) \geq f([aB]) + f(b^2) \geq f([aB])
				\]
				as before. This proves unoriented connected smoothing.
				
			\end{itemize}
\end{itemize}

						This finishes the proof that disconnected smoothing implies
		connected smoothing.
\end{proof}	
		
\begin{proof}[Proof of Direction 2 of Proposition~\ref{prop:cont_conn_disc}]
		
		We now prove the second part, namely, if $f \colon \Curves^+(S) \to \mathbb{R}$ satisfies symmetry, connected smoothing, homogeneity and continuity, then $f$ satisfies disconnected smoothing. The proof is very similar to the first part.
		
		\begin{itemize}
			\item \textbf{Part one:} Oriented disconnected smoothing $[a][b] \reducesto
			[ab]$.
			
			 Namely, given $a, b \in \pi_1(S)$ with $(\ora{a},\ora{b})$ crossing, we prove, by induction on $n$, the following inequality
			\begin{equation}
				f([a^{2n}b^{2n}]) \geq 2nf([ab]).
				\label{eq:inequality_smoothing}
			\end{equation}
			
			The base case follows by noting that if $(\ora{a},\ora{b})$ R-cross, then $(\ora{ba},\ora{ab})$ are R-parallel, and then by connected smoothing
			\[
			f([a^2b^2])=f([ab^2a]) \geq 2f([ab]).
			\]
			For the induction step, note that
			\begin{align*}
				f([a^{2n}b^{2n}]) &= f([a^2(a^{2(n-1)}b^{2(n-1)})b^2])\\
				&= f([b^2a^2(a^{2(n-1)}b^{2(n-1)})]) \\
				&\geq f([b^2a^2]) + f([a^{2(n-1)}b^{2(n-1)})]) \\
				&\geq 2f([ab]) + 2(n-1)f([ab]) = 2nf([ab]).
			\end{align*}
			To get the first inequality we used oriented connected smoothing and that $b^2a^2$ and $a^{2(n-1)}b^{2(n-1)}$ are parallel, which follows from the fact that $a$ and $b$ are crossing, and Lemma~\ref{lem:anbn_vs_bkak_parallel}. For the second inequality we used the induction step twice.
			Finally, using Lemma~\ref{lem:gnhn_to_gsumh} as above, we get
			\[
			f([a] [b]) \geq f([ab])
			\]
			which proves oriented disconnected smoothing.
	
			
			\item \textbf{Part two:} Unoriented disconnected smoothing $[a][b]\reducesto [aB]$.
			
			Unoriented disconnected smoothing is proved in the same way, replacing
			$b$ by~$B$.
	\end{itemize}
\end{proof}

As an immediate corollary, we get the following.
\DPTtodo{Do we really need two corollaries that are basically
  replacing ``convex union'' to ``additive''? Why are they separated
  so far in the text?}
  \DMGtodo{we don't need it, I deleted it and referred to the proposition in the relevant places}

Finally, we show that power smoothing and oriented disconnected smoothing imply stability.

\begin{proposition}
  \DPTtodo{Why are we spelling out the conditions here and not in
    other statements in this section?}
    \DMGtodo{removed the spelling-out}
Let $f \colon \Curves^+(S) \to \mathbb{R}$ be a curve functional satisfying:
\begin{enumerate}
\item additivity
\item sub-stability
\item oriented disconnected smoothing.
\end{enumerate}
\DPTtodo{Note ``stability'' as previously defined applies to curve
  functionals, not functions on $\pi_1$}
  \DMGtodo{stated it in terms of a curve functional}
Then $f$ satisfies stability: $f(C^n) = nf(C)$.
\label{prop:power_disconnected_imply_stable}
\end{proposition}
\begin{proof}
  Let $a \in \pi_1(S)$ be an arbitary element.
\DPTtodo{Reorganized the proof to put the roadmap first}
By hypothesis~(2)
\DPTtodo{We have not used ``sub-stable'' before. It's a fine term that
  maybe we should introduce earlier} \DMGtodo{I'm okay with calling it sub-stable. Also, I fixed power smoothing to mean: $f([a^{n+m}]) \geq f([a^n][a^m])$. Sub-stable means $f([a^{n+m}]) \geq f([a^n])+ f([a^m])$  which is really what we need here. Of course, power smoothing + additive implies sub-stable. I added definition of sub-stable after definition of stable, and also added this remark}
(sub-stable), we have $f(a^{n+m}) \geq f(a^n) + f(a^m)$.
By Fekete's Lemma~\ref{lem:fekete}, the limit
\[
F(a)\coloneqq\lim_{n^+} \frac{f(a^n)}{n}
\]
exists, and $f(a) \leq F(a)$. We must show $F(a) \leq f(a)$.
  
  Find $b \in
\pi_1(S)$ so that $(\ora{a},\ora{b})$ R-cross, which we can find by
\DPTtodo{We already showed we could find $b$ like this explicitly and
  more precisely, cite that}
  \DMGtodo{yes, in the proof of Lemma~\ref{lem:if_unbounded_then_reducible}, but it is much later in the paper. Should we add a forward reference or be precise here and add a backward reference there? There is also another instance before, in the proof of Proposition~\ref{prop:graph_collapse}, but we are less precise there, and also the crossing elements are assumed to satisfy further properties. I did not change anything.}
density of endpoints of axes of hyperbolic elements in $G(S)$.
Then, applying oriented disconnected smoothing to the R-crossing pair $(\ora{a^n}, \ora{b})$, we get
\begin{equation}
f(a^n) + f(b) \geq f(a^nb).
\label{eq:cross1_eq}
\end{equation}

We next show that
\begin{equation}
nf(a) + f(b) \geq f(a^nb)
\label{eq:cross2_eq}
\end{equation}
for $n\geq 1$, by induction on $n$.
Indeed, the base case is true by oriented connected smoothing, and the
induction follows from
\[
nf(a) + f(b) \geq f(a^{n-1}b) + f(a) \geq f(a^nb)
\]
where in the first step we used the induction step and in the second we used the fact that the pair $(\ora{a},\ora{a^{n-1}b})$ is R-crossing by Lemma~\ref{lem:rhs_prod}.

Finally, by Lemma~\ref{lem:rhs_prod}, the pair $(\ora{B},\ora{a^nb})$
is R-crossing for all $n \geq 1$, and hence, by oriented disconnected
smoothing,
\begin{equation}
f(B) + f(a^nb) \geq f(a^n).
\label{eq:cross3_eq}
\end{equation}

Combining Eqs.~\eqref{eq:cross1_eq}, and~\eqref{eq:cross3_eq},
dividing by $n$ and taking the limit when $n$ goes to infinity, we get
$F(a)=\lim_{n^+} \frac{f(a^nb)}{n}$. Using Eq.~\eqref{eq:cross2_eq}
we get $f(a) \geq F(a)$. Thus, $F(a)=f(a)$, and hence $f$ satisfies stability since $F$ clearly does.
\end{proof}

\begin{remark}
If we modify the statements of
Propositions~\ref{prop:oriented_smoothing_unoriented}
and~\ref{prop:cont_conn_disc} to refer to $C$-quasi-smoothing (for fixed
$C>0$), then the proofs \emph{do not} go through. For example, in the
induction from connected to disconnected
oriented smoothing, inequality~\ref{eq:inequality_smoothing} would
have to be replaced by
\[
f([a^{2n}b^{2n}]) \geq 2nf([ab]) - C n.
\]
Following the same chain of inequalities that prove the induction does
not yield the desired inequality, since one accumulates one extra $-C$
factor. We suspect the modified statement for quasi-smoothing is not
true.
\end{remark}

%%% Local Variables:
%%% mode: LaTeX
%%% TeX-master: "Intersections"
%%% End:
